
\subsection{Meeting of 27.01.26}
\subsubsection{Since Last Time}
\begin{itemize}
    \item Setup of Julia environment    
    \item Basic implementation of anistropic GNN which was used in the Difusco paper
    \item Unit test for the anisotropic GNN layer - compare with manual calculations.
\end{itemize}
\subsubsection{Questions}
\begin{itemize}
    \item Next step for me: Try to define loss for MISP and then see whether we see some results. Is this ok?
    \item Should I use the instance structure from the GR implementation?
    \item Is there a MISP dataset, or should I just generate one? Got this Julia package that does that anyway.
\end{itemize}
\subsubsection{Unstructured Notes}
\begin{itemize}
    \item Graph types: KRegularRandomgraphs für MISP might make sense. 
    \item Another option barabasi-albert graphs - but not checked yet.
    \item RI graphs also another classic option. (Be careful not to dense).
    \item GeometrixFlux.jl package nochmal anschauen, da gibts eine GNN implementierung, this might al include the loss.
    \item Reuse GNNs from Libs.
    \item GNN.jl lib kann GNN auch ersetzen evtl. 
\end{itemize}